\thispagestyle{empty}
    \begin{minipage}{0.22\linewidth}
    \begin{tcolorbox}[colback=white, title=, width=1\linewidth]	
        BTS 1ère Année \\
        Spécialité CIEL 
    \end{tcolorbox}
\end{minipage}\hfill
\begin{minipage}{0.55\linewidth}
    \begin{tcolorbox}[colback=white, title=, width=1\linewidth]	
        \begin{center}
            Lycée Victor Hugo~--~COLOMIERS \\
            CCF Mathématiques. Epreuve E3~--~Sous épreuve U31 
        \end{center}
        
    \end{tcolorbox}
\end{minipage}\hfill
\begin{minipage}{0.23\linewidth}
    \begin{tcolorbox}[colback=white, title=, width=1\linewidth]	
        \begin{center}
            Lundi 04 Mai 2026 \\
    Durée 55 min
        \end{center}
        
    \end{tcolorbox}
\end{minipage}\hfill

\vspace{2cm}

\NomPrenom{}

\vspace{2cm}

\begin{center}
    
\emph{Un ordinateur doté d'un tableur et du logiciel Géogébra est mis à la disposition du candidat qui pourra également se servir de sa calculette scientifique.}

\bigskip

\emph{Aucun document n'est autorisé.}

\bigskip

\emph{L'usage du téléphone est interdit.}

\bigskip

\emph{La clarté du raisonnement et la qualité de la rédaction entreront pour une part importante dans l'appréciation des copies.}

\end{center}

\bigskip

\vspace{2cm}

Ce sujet comporte~\pageref{LastPage} pages numérotées de 1 à~\pageref{LastPage} et est constitué de 3 parties indépendantes ainsi que de deux appels professeur.

\vspace{2cm}

\section*{Évaluation des compétences \hfill \makebox[2cm]{\dotfill} / 10}

\renewcommand{\arraystretch}{2}

\begin{tabularx}{\textwidth}{|X|c|c|}
\hline
\rowcolor{gray!20} \textbf{Compétences évaluées} & \textbf{Questions} & \textbf{Notes} \\ \hline
\textbf{S'informer}: Rechercher, extraire et organiser l'information. & Q1, Q5, Q9, Q11, Q18& \makebox[1cm]{\dotfill} / 1\\ \hline
\textbf{Chercher}: Proposer une méthode, expérimenter, tester.& Q8, Q11, Q14, Q22&\makebox[1cm]{\dotfill} / 2\\ \hline
\textbf{Modéliser}: Traduire un problème en langage mathématique. &Q1, Q3, Q17, Q21&\makebox[1cm]{\dotfill} / 1\\ \hline
\textbf{Raisonner}: Déduire, induire, justifier ou démontrer. &Q7, Q10, Q13, Q19&\makebox[1cm]{\dotfill} / 2\\ \hline
\textbf{Calculer}: Calculer, illustrer, programmer, mettre en œuvre. & Q2, Q3, Q4, Q6, Q9, Q12, Q15&\makebox[1cm]{\dotfill} / 2\\ \hline
\textbf{Communiquer}: Rendre compte d'une démarche à l'oral ou à l'écrit. & Q4, Q8, Q11, Q13, Q16&\makebox[1cm]{\dotfill} / 2\\ \hline
\end{tabularx}
\renewcommand{\arraystretch}{1}

\vfill

\newpage